\ifx\danslelivre\undefined
\documentclass[10pt]{../fichiers-configuration-latex/smfart}
\usepackage{ifthen}
\newboolean{english}
\newboolean{francais}
\setboolean{english}{false}
\setboolean{francais}{true}
\usepackage{palatino,euler}
\input{../fichiers-configuration-latex/paquets}
\selectlanguage{francais}
\input{../fichiers-configuration-latex/lemme-proposition-theoreme}
\input{../fichiers-configuration-latex/numerotation}
\input{../fichiers-configuration-latex/macros}
\input{../fichiers-configuration-latex/ucs_manquants}
\input{../.cv}

\usepackage{tikz}
\usetikzlibrary{matrix}
\usetikzlibrary{calc}
\usetikzlibrary{decorations.markings}

\externaldocument{../00-introduction/introduction}
\externaldocument{../01-quasi-excellents/excellent}
\externaldocument{../03-approximation/approximation}
\externaldocument{../04-Cohen-Gabber/Cohen-Gabber}
\externaldocument{../05-algebrisation-partielle/algebrisation}
\externaldocument{../06-log-lissite/log}
\externaldocument{../07-uniformisation-locale/uniformisation-faible}
\externaldocument{../08-modification-theorem/modification}
\externaldocument{../09-uniformisation-premiere-a-l/uniformisation-locale-premiere-a-l}
\externaldocument{../11-produits-orientes/produits-orientes}
\externaldocument{../12-descente-cohomologique-orientee/descente-orientee}
\externaldocument{../13-theoreme-de-finitude/finitude}
\externaldocument{../14-fonctions-dimension/fonctions-dimension}
\externaldocument{../15-Lefschetz-affine/Lefschetz-affine}
\externaldocument{../16-classes-de-Chern/gysin}
\externaldocument{../17-complexes-dualisants/dualite}
\externaldocument{../18A-cohomological-dimension/dimensionA}
\externaldocument{../19-exemples-contre-exemples/contre-exemples}
\externaldocument{../20-rigidite/rigidite}
\externaldocument{../21-finitude-non-abelienne/finitude-non-abelienne}

\begin{document}

\begin{center}
XVII-B. Dimension cohomologique : raffinements et compléments\\
Fabrice Orgogozo
\end{center}

\version
\tableofcontents

\else
\setboolean{english}{false}
\setboolean{francais}{true}
\selectlanguage{francais}
\chapter{Dimension cohomologique : raffinements et compléments}
\addtocontents{toc}{par Fabrice Orgogozo}
\begin{center}
Fabrice Orgogozo
\end{center}
\fi

\newcommand{\initial}{\mathop{\mathrm{in}}}

Le premier objectif de cette seconde partie de l'exposé est de
démontrer que pour tout nombre premier $ℓ$ il existe un ouvert
\emph{affine} d'un schéma nœthérien strictement hensélien régulier de
dimension $2$ dont la $ℓ$-dimension cohomologique est égale à $3$.
Outre les ingrédients cohomologiques — pureté, morphisme de Gysin et
comparaison à la complétion —, on utilise une construction dont le
principe est dû à Nagata : utilisant des « dilatations formelles »
construit un schéma nœthérien strictement hensélien $X$ de
dimension $2$, de complété $\chap{X}$ \emph{régulier de
dimension $2$}, et une courbe \emph{irréductible} $C$ dans $X$
devenant le $ℓ$-ième multiple d'un diviseur régulier dans $\chap{X}$.
Cette construction est ensuite étendue au cas, plus délicat, de la
dimension supérieure. À partir de là, on construit
aisément des schémas dont l'existence a été annoncée
dans la première partie (\refext{XVIII-A}{18A-remark-d-2d-1}).
Pour vérifier que leur dimension cohomologique est bien celle
attendue, on fait appel à une majoration
assez générale établie sans hypothèse d'excellence.
Enfin, on termine par une minoration de la dimension
cohomologique d'un ouvert (non nécessairement
affine) du spectre épointé d'un schéma nœthérien intègre
strictement local.

\section{Préliminaires}

\subsection{Dilatations formelles}
\label{18B-dilatations-formelles}

Soient $R$ un anneau, $π$ un élément non diviseur de $0$
et $f$ un élément de la complétion $π$-adique $\chap{R}$
de $R$. Pour tout $n ≥ 0$, choisissons un $f_n ∈ R$ tel que
$f ≡ f_n$ modulo $π^n$.

\begin{définition2}
On note $𝒟𝒾𝓁^f_π R$ la sous-$R$-algèbre de $R[π^{-1},F]$, où $F$ est une
indéterminée, colimite des $R$-algèbres $R[\frac{F-f_n}{π^n}]$.
\end{définition2}

On notera également $F$ l'image de cette variable dans
$𝒟𝒾𝓁^f_π R$.

\begin{remarque2}
Notons que les $R$-algèbres considérées sont toutes isomorphes
à une algèbre de polynômes en une variable sur $R$.
\end{remarque2}

\subsubsection{}On vérifie immédiatement les faits suivants :

\begin{enumerate}
\item la construction ne dépend pas des choix des $f_n$,
et ne dépend de l'élément $π$ qu'à travers l'idéal qu'il
engendre ;
\item les morphismes $R →  𝒟𝒾𝓁_π^f R$
et $𝒟𝒾𝓁_π^f R → \chap{R}$, $\frac{F-f_n}{π^n}↦ \frac{f-f_n}{π^n}$,
induisent des \emph{isomorphismes} sur
la complétion $π$-adique.
\end{enumerate}

\subsection{Platitude et nœthérianité}

\subsubsection{}Rappelons que si un morphisme $A → B$ est fidèlement plat,
$A$ est nœthérien si $B$ l'est (\ega{IV₂}{2.2.14}).
Pour vérifier la platitude, il est parfois commode d'utiliser
le critère suivant (\cite[II.1.4.2.1]{Platification@Raynaud-Gruson}).

\begin{proposition2}
\label{18B-critère-platitude}
Soient $M$ un $R$-module, et $π ∈ R$. On suppose que $π$ n'est diviseur de zéro
ni dans $R$ ni dans $M$.  Alors, $M$ est \emph{plat} sur $R$ si et seulement si
$M/π$ l'est sur $R/π$ et $M[π^{-1}]$ l'est sur $R[π^{-1}]$.
\end{proposition2}

\begin{remarque2}
Pour démontrer la nœthérianité des anneaux considérés
ci-après, on pourrait également utiliser le critère
de Cohen rappelé en \refext{XIX}{19-copremier},
en vérifiant notamment que les idéaux de hauteur $1$
sont principaux.
\end{remarque2}

\subsection{Gonflements}
\label{18B-gonflement}
\subsubsection{}Soit $A$ un anneau local nœthérien, d'idéal maximal $𝔪$.
Suivant \BourbakiAC{IX, appendice, §2}, on note $A]t[$
et on appelle « gonflement (élémentaire) »
de $A$ le localisé de l'anneau de polynômes $A[t]$ en l'idéal premier $𝔪A[t]$.
C'est un anneau local nœthérien.
(Il est noté $A(t)$, par analogie avec les
fractions rationnelles, dans \cite[p. 17—18]{Nag62} ;
voir aussi  \cite[p. 138]{KKK@Matsumura}.)
Plus généralement, on peut considérer un ensemble arbitraire
de variables $t_{e ∈ E}$ et définir l'anneau $A]t_e, e ∈ E[$,
localisé de $A[t_e, e ∈ E]$ en l'idéal premier engendré
par $𝔪$. Rappelons le fait suivant (\BourbakiAC{IX, appendice, prop. 2
et corollaire}).

\begin{proposition2}
L'anneau $A]t_e, e ∈ E[$ est local nœthérien de même dimension
que $A$.
\end{proposition2}

\subsubsection{}Notons que le cas d'un nombre fini de variables est
très élémentaire et que le cas général
résulte du lemme \ega{0$_{III}$}{10.3.1.3},
reproduit en \refext{XIX}{19-ega-noeth},
par passage à la (co)limite.
Pour une autre démonstration, voir également
\BourbakiAC{III.§5, exercice 7}.

\subsubsection{}
Notons que si le corps résiduel de $A$ est $κ$, celui
de $A]t_e, e ∈ E[$ est canoniquement isomorphe à son
extension transcendante pure $κ(t_e, e ∈ E)$.
De plus, on montre que si $F$ est un sous-ensemble de $E$,
le morphisme $A]t_e, e ∈ F[ → A]t_e, e ∈ E[$ est \emph{fidèlement plat}.
(Voir \BourbakiAC{IX, appendice, prop. 2} pour une démonstration
dans le cas  $F=∅$, auquel on se ramène immédiatement.)

\section{Construction de Nagata en dimension $2$, application
cohomologique}
\label{18B-Nagata-Gabber-dimension-2}

\subsection{Dilatation relativement à une série transcendante}

\subsubsection{}Soit $W$ un anneau de valuation discrète,
d'idéal maximal $𝔪_W=(π)$, corps résiduel $k$, corps des
fractions $K$ et de complété $\chap{W}$. On note
$\chap{K}$ le corps des fractions de $\chap{W}$.
Supposons qu'il existe un élément $φ∈ \chap{W}$ \emph{transcendant}
sur $K$. C'est le cas si $W$ est dénombrable ou, par exemple,
lorsque $W=k[[t]]$ auquel cas $φ=∑_n t^{n!}$ convient.
Par translation, on peut supposer que $φ$ appartient à l'idéal maximal de $\chap{W}$.

\subsubsection{}
Fixons un entier $ℓ ≥ 1$ et considérons l'élément $f=(y - φ)^ℓ$ du complété $π$-adique
$\chap{W[y]}$ de $W[y]$. Notons que ce complété s'injecte dans la complétion totale
$\chap{W}[[y]]$ et que $f$ appartient au sous-anneau $\chap{W}[y]$
de $\chap{W[y]}$ et $\chap{W}[[y]]$. En conséquence, le morphisme
canonique $𝒟𝒾𝓁_{π}^{f} W[y] → \chap{W[y]}$ se factorise
en un morphisme $𝒟𝒾𝓁_{π}^{f} W[y] → \chap{W}[y]$.

\begin{proposition2}
Le morphisme $𝒟𝒾𝓁_{π}^{f} W[y] → \chap{W}[y]$
est \emph{plat}.
\end{proposition2}

\begin{démo}
Notons pour simplifier $𝒟$ la source de ce morphisme.
D'après le critère de platitude rappelé ci-dessus,
il suffit de montrer que le morphisme
$𝒟[π^{-1}] → \chap{W}[y][π^{-1}]$ est plat
car $𝒟 → \chap{W}[y]$ induit visiblement un isomorphisme
modulo $π$. Or, lorsque l'on inverse $π$, les $W[y]$-algèbres
dont $𝒟$ est par définition la colimite sont
toutes isomorphes à un anneau de polynômes $K[y,F]$ où,
rappelons-le, $K$ est le corps des fractions $W[π^{-1}]$
de $W$. Par conséquent, il nous faut montrer
que le morphisme $K[y,F] → \chap{W}[y][π^{-1}]$,
composé des morphismes

\begin{minipage}[c]{.46\linewidth}
\begin{center}
\begin{equation*}
\begin{cases}
K[y,F] → K[y, F′]\\
F ↦  {F′}^ℓ
\end{cases}
\end{equation*}
\end{center}
\end{minipage} \hfill
\begin{minipage}[c]{.46\linewidth}
\begin{center}
\begin{equation*}
\begin{cases}
K[y, F′] → \chap{K}[y] \\
F′↦ y-φ
\end{cases}
\end{equation*}
\end{center}
\end{minipage}
La platitude du premier est évidente.
Pour le second, on se ramène par translation et changement
de base à montrer que le morphisme $K[F′] → \chap{K}$, $F′↦ φ$,
est plat. Il se factorise en le composé du passage aux fractions
$K[F′] → K(F′)$ avec l'\emph{injection} $K(F′) ↪ \chap{K}$
déduite de $φ$. Chacun de ces morphismes est plat.
\end{démo}

\begin{remarque2}
La construction de l'anneau de dilatation $𝒟𝒾𝓁_{π}^{f} W[y]$
est inspirée de celle de \cite{example@Nagata}, qui considère
le cas $ℓ=2$. (Voir aussi \cite[appendice, E4.1]{Nag62}
et \cite{Noetherian@HRW}.)
\end{remarque2}

\subsection{Le diviseur $C=V(F)$}

\subsubsection{} On conserve les notations précédentes.
Soit $A$ le localisé de l'anneau de dilatation $𝒟$ en l'idéal
premier image inverse de l'idéal $(π,y) ⊆ \chap{W}[y]$.
C'est un anneau nœthérien par fidèle platitude, de corps
résiduel $k$.

\begin{lemme2}
L'anneau $A$ satisfait les propriétés suivantes :
\begin{enumerate}
\item il est régulier et la suite $(π,F)$ est régulière ;
\item son quotient $A/F$ est intègre ;
\item l'intersection schématique du fermé $V(F)$ avec le spectre épointé
de $A$ est un schéma régulier.
\end{enumerate}
\end{lemme2}

\begin{démo} 
(i) Il suffit d'établir les deux énoncés pour $𝒟$.
Pour se faire, on peut compléter $π$-adiquement
(cf. par exemple \BourbakiAC{X, §4, nº2, cor. 3}).
Le complété de $𝒟$ est isomorphe à $\chap{W[y]}$
de sorte que la régularité de l'anneau et de
la suite $(π,F)$ — c'est-à-dire l'injectivité de
la multiplication par $π$, et par $F$ modulo $π$— sont évidents.
(ii) Si on inverse $π$, l'anneau $A$ devient une localisation de l'anneau de polynômes $K[y,F]$.
La restriction du diviseur à cet ouvert est intègre.
La multiplication par $π$ dans le quotient $A/F$
étant injective d'après (i), l'intégrité de $A/F$
résulte de celle de $A[π^{-1}]/F$.
Le quotient est non nul car $F ∈ 𝔪_A$.
(iii) Sur l'ouvert complémentaire de $V(π)$, l'élément $F$
est une indéterminée de sorte que le résultat est clair.
D'autre part, l'intersection du complémentaire de $V(y)$
avec le diviseur est contenu dans le complémentaire de $V(π)$
car l'équation est — dans le complété $π$-adique — de la forme $(y-φ)^ℓ$, avec $φ
∈ (π)$. Ceci suffit pour conclure.
\end{démo}

\subsubsection{}Notons le fait suivant, trivial mais crucial :
par construction, le diviseur $V(F)$ devient
$V((y - φ)^ℓ)$ dans le complété $\chap{A}$ de $A$
relativement à son idéal maximal.

\subsection{Hensélisation}
Pour simplifier les notations, on suppose
dorénavant le corps $k$ séparablement clos.

\subsubsection{}
Soient $A^\hs$ le hensélisé de $A$ en son point fermé,
$\chap{A}$ le complété de $A$ (ainsi que de $A^\hs$),
et notons $X=\Spec(A^\hs)$ et $\chap{X}=\Spec(\chap{A})$
leurs spectres respectifs, ainsi que $⋆$ et $\chap{⋆}$ les points fermés.
Comme $\Spec(A)$, le schéma $X$ est intègre.
De plus, le diviseur $C$ d'équation $F=0$ dans $X$
est réduit, cette propriété étant
également conservée par hensélisation.

\subsubsection{}
Vérifions que le diviseur $C$ est \emph{irréductible}.
D'après le théorème de comparaison de Elkik,
le morphisme $π₀(\chap{C}-\chap{⋆}) → π₀(C - ⋆)$ est une
bijection. Or, $\chap{C}-\chap{⋆}$ est connexe :
dans un anneau local régulier $B$, le spectre
$\Spec(B/g^ℓ)$ est irréductible pour tout $g ∈ 𝔪_B-𝔪_B²$.
Ainsi, l'ouvert $C - ⋆$ de $C$ est connexe et, finalement,
$C$ est irréductible.

\subsection{Application cohomologique}
\emph{On suppose dorénavant l'entier $ℓ$ inversible sur $X$.}

\subsubsection{}
Notons pour simplifier $G$ l'élément $(y-φ)^ℓ$ de $\chap{A}$,
de sorte que — par construction — on a l'égalité $F=G^ℓ$ dans $\chap{A}$.
Notons $U$ l'ouvert affine $X-C$ du schéma strictement local $X$.
Soit $j$ l'immersion ouverte de $U$ dans le spectre
épointé $X-⋆$ et $i$ l'immersion fermée $C-⋆ ↪ X-⋆$.

\subsubsection{}Le triangle
\[
i_* i^! → \mathrm{Id} → 𝖱 j_* j^*
\]
sur $X-⋆$ induit la suite exacte
\[
H³_{C-⋆}(X-⋆,𝐙/ℓ) → H³(X-⋆,𝐙/ℓ) → H³(U,𝐙/ℓ) →
H⁴_{C-⋆}(X-⋆,𝐙/ℓ)=H²(C-⋆,i^! 𝐙/ℓ[2]).
\]
Par pureté (\refext{XVI}{16-theoreme-purete-absolue}), le faisceau de cohomologie locale $i^! 𝐙/ℓ$ est
constant concentré en degré $2$. Or le groupe de cohomologie $H²(C-⋆,𝐙/ℓ)$ est
nul : la cohomologie du corps des fractions d'un anneau $B$
strictement local intègre de dimension $1$ est nul en
degré $≥2$. (On se ramène au cas bien connu d'un anneau de valuation
discrète en observant que le normalisé
de $B$ dans son corps des fractions est un
anneau nœthérien, de Dedekind, et strictement local
car colimite locale d'anneaux strictement locaux.)
Il en résulte que la flèche de restriction
$H³(X-⋆,𝐙/ℓ) → H³(U,𝐙/ℓ)$ est surjective ; nous allons
voir que c'est un isomorphisme.

\subsubsection{}Comme rappelé ci-dessus, le
morphisme $H³_{C-⋆}(X-⋆,𝐙/ℓ) → H³(X-⋆,𝐙/ℓ)$
s'identifie, par pureté, au morphisme
de Gysin $\mathrm{Gys}(f):H¹(C-⋆,𝐙/ℓ) → H³(X-⋆,𝐙/ℓ)$.
Il résulte de la commutativité du diagramme
\begin{center}
 \begin{tikzpicture}[auto]
  \matrix(diag)[matrix of math nodes,column sep=8em,row sep=5ex]
  {|(D)|  H¹(C-⋆,𝐙/ℓ)   & |(X)| H³(X-⋆,𝐙/ℓ)  \\
   |(Dc)|  H¹(\chap{C}-⋆,𝐙/ℓ) & |(Xc)| H³(\chap{X}-⋆,𝐙/ℓ),  \\};
  \draw[->] (D) -- (Dc);
  \draw[->] (D) -- node{$\mathrm{Gys}(F)$} (X);
  \draw[->] (Dc) -- node{$\mathrm{Gys}(F)=\mathrm{Gys}(G^ℓ)$} (Xc);
  \draw[->] (X) -- (Xc);
 \end{tikzpicture}
\end{center}
de l'égalité $\mathrm{Gys}(F)=ℓ ⋅ \mathrm{Gys}(G)$
et enfin du fait que les flèches verticales sont
des isomorphismes (comparaison à la complétion,
\cite[6.6.4]{TUBULAR@Fujiwara}) que le morphisme $\mathrm{Gys}(F)$ est nul.
(La commutativité du diagramme résulte par exemple
de la définition \refext{XVI}{16-definition-cl-i}
et de \refext{XVI}{16-proposition-classes-generalisees-cb}.)
Ainsi, le morphisme de restriction induit un
isomorphisme
\[
H³(X-⋆,𝐙/ℓ) ⥲ H³(U,𝐙/ℓ).
\]
Or le terme de gauche est non nul, à nouveau par pureté.
Le schéma affine $U$ est donc de dimension cohomologique $>2$. CQFD.

\section{Séries formelles de Gabber, application cohomologique}
\label{18B-Gabber-Nagata-dimension-supérieure}

On étend la construction précédente en dimension
arbitraire $≥ 2$.

\subsection{Une série formelle et sa décomposition}

\subsubsection{}
\label{18B-décomposition-série-formelle}
Soit $A$ un anneau commutatif.
Rappelons que l'application $A$-linéaire
\[
∑_{i=1}^n A[x_{j ≠ i}][[x_i]] → A[[x₁,…,x_n]]
\]
somme des injections canoniques est \emph{surjective} :
si $G ∈  A[[x₁,…,x_n]]$, on peut par exemple regrouper pour chaque $i ∈ [1,n]$
les termes $ax₁^{β₁}\cdots x_n^{β_n} ∈ A[[x₁,…,x_n]]$ de $G$
pour lesquels $β_i=\max_{j ∈ [1,n]} β_j$ et $β₁,…,β_{i-1}<β_i$.
(Cette dernière condition n'est là que pour définir $i$
de façon non ambiguë ; tout autre choix conviendrait.)
La somme $g_i$ de ces termes appartient à $A[x_{j ≠ i}][[x_i]]$, et
$G=g₁+\cdots+g_n$.

\subsubsection{}On fixe maintenant deux entiers non nuls $n$ et $ℓ$
et on considère
\[
S=\big(y + ∑_{i=1}^n ∑_{α=1}^{∞} t_{i α} x_i^α\big)^ℓ ∈ 𝐙[y,t_{i
∈ [1,n], α ≥ 1}][[x₁,…,x_n]].
\]
Il résulte de l'observation précédente que l'on peut écrire cette
série sous la forme
\[
y^ℓ + f₁+\cdots+f_n
\]
où chaque $f_i$ est une série formelle en $x_i$, à
coefficients polynomiaux en les autres variables.

\subsubsection{}
\label{18B-définition-fi}
Afin que la proposition de platitude ci-dessous soit vraie,
on procède de façon légèrement différente pour définir les
séries formelles $f_i ∈ 𝐙[y,t_{jα},x_{k ≠ i}][[x_i]]$ telles
que $S-y^ℓ=∑_{i=1}^n f_i$. Écrivons
\[
S=y^ℓ + ∑_{α=1}^∞ (∑_{i=1}^n t_{i α}x_i^α)^ℓ +
\big(\text{élément de degré }<ℓ\text{ en les }t_{jβ}\big).
\]
Soient $i ∈ [1,n]$ et $α ≥ 1$ des indices.
On considère les termes $ax₁^{β₁}\cdots x_n^{β_n}$ de $(∑_{i=1}^n t_{i α}x_i^α)^ℓ$
pour lesquels $i$ est le plus grand indice tel que $β_i ≠ 0$,
c'est-à-dire les termes de $(∑_{i=1}^n t_{i α}x_i^α)^ℓ$
qui sont dans $𝐙[t_{1 α},…,t_{n α},x₁,…,x_i]$ mais pas
dans $𝐙[t_{1 α},…,t_{n α},x₁,…,x_{i-1}]$.
À $i$ fixé, la somme sur $α$ de ces termes
est un élément $f_{i,=ℓ}$ de $𝐙[t_{jβ},x_{k ≠ i}][[x_i]]$.
Par construction, on a l'égalité
$ ∑_{α=1}^∞ (∑_{i=1}^n t_{i α}x_i^α)^ℓ=∑_{i=1}^n f_{i,=ℓ}$.
Enfin, on décompose le terme restant,
$S-y^ℓ-∑_{i=1}^n f_{i,=ℓ}$, en une somme $∑_{i=1}^n f_{i,<ℓ}$
où chaque $f_{i,<ℓ}$ appartient à $𝐙[y,t_{jβ},x_{k ≠ i}][[x_i]]$,
en procédant par exemple comme en \ref{18B-décomposition-série-formelle}.
On pose alors $f_i=f_{i,=ℓ}+f_{i,<ℓ}$ ; chacun de ces
éléments est de degré total en les $t_{j β}$ inférieur ou égal à $ℓ$,
et appartient donc également à $𝐙[y,t_{j,β ≠ α}][[x₁,…,x_n]][t_{1
α},…,t_{n α}]$ pour chaque $α ≥ 1$.

\begin{proposition2}
\label{18B-platitude-série-formelle-universelle}
Fixons $α ≥ 1$. Notons $T_i=t_{i α}$ pour chaque $i ∈ [1,n]$
et $R_α$ l'anneau $𝐙[y,t_{j,β ≠ α}][[x₁,…,x_n]]$.
Le morphisme
\[
R_α[F₁,…,F_n] → R_α[T₁,…,T_n]
\]
\[
F_i↦ f_i
\]
est \emph{libre}, donc \emph{plat}, au-dessus de l'ouvert
$x₁\cdots x_n ≠ 0$ de $\Spec(R)$.
\end{proposition2}

\begin{démo}
Par construction, chaque $f_i$ est une somme $f_{i,=ℓ} + f_{i,<ℓ}$, où
$f_{i,=ℓ}$ (resp. $f_{i,<ℓ}$) est un polynôme dans $R_α[T₁,…,T_n]$ de
degré total égal (resp. strictement inférieur) à $ℓ$.  De plus,
$f_{i,=ℓ}$ est, comme polynôme en $T_i$, de la forme $x_i^{αℓ} T_i^ℓ +
∑_{m<ℓ} c_m T_i^m$ où les $c_m$ appartiennent à $R_α[T_{j<i}]$.
Munissons les monômes de $R_α[T₁,…,T_n]$ de l'ordre lexicographique
gradué suivant : 
$T₁^{d₁} \cdots T_n^{d_n} \mathrel{\preceq} T₁^{d₁′} \cdots T_d^{d_n′}$
si et seulement si
$\sum d_i < \sum d′_i$ ou $\sum d_i = \sum
d′_i$ et $d_i < d'_i$ pour le plus grand $i$ tel que
$d_i \neq d'_i$.
Il est clair que le \emph{terme de tête} $\initial_{\preceq}(g)$
pour cet ordre d'un polynôme $g=f₁^{q₁}\cdots f_n^{q_n}$
en les $f_i$ est $T₁^{q₁ ℓ}\cdots T_n^{q_n ℓ}$, à multiplication près
par un monôme en les $x_i$.
Il en résulte immédiatement qu'en inversant les $x_i$,
l'anneau $R_α[T₁,…,T_n]$ est libre sur $R_α[F₁,…,F_n]$
de base les monômes $T₁^{r₁}\cdots T_d^{r_d}$,
avec $0 ≤ r_i < ℓ$.
\end{démo}

\subsubsection{}Soit $A$ un anneau local nœthérien régulier
de dimension $n$ dont on note $x₁,…,x_n$ un système
régulier de paramètres. On pourra penser par
exemple au localisé $k[x₁,…,x_n]_{(x₁,…,x_n)}$ d'un anneau
de polynômes sur un corps.
On considère le gonflement $A]\sous{t}[$ défini
en \ref{18B-gonflement}, où $\sous{t}$
est l'ensemble des variables $\{t_{i α}: i ∈ [1,…,n], α
∈ 𝐍_{≥1}\}$.

\subsection{Construction d'un anneau local régulier pathologique}
Soit $i ∈ [1,…,n]$. Notons encore $f_i$
l'image de la série formelle à coefficients entiers
considérée en \ref{18B-définition-fi}
dans le complété $x_i$-adique de $A]\sous{t}[ [y]$,
et $𝒟_i$ l'anneau de dilatation
$𝒟𝒾𝓁_{x_i}^{f_i} A]\sous{t}[ [y]$.
Le produit tensoriel $𝒫$ de ces $A]\sous{t}[ [y]$-algèbres s'envoie naturellement
dans le complété $\chap{A]\sous{t}[} [[y]]$,
où la première complétion est faite relativement à l'idéal
maximal $(x₁,…,x_n)$ de $A]\sous{t}[$.
On note $𝒟$ le localisé de $𝒫$ en l'image de l'idéal maximal $(x₁,…,x_n,y)$ de
$\Spec(\chap{A]\sous{t}[} [[y]])$.

\begin{proposition2}
Le morphisme $𝒟 → \chap{A]\sous{t}[} [[y]]$ est fidèlement \emph{plat}.
\end{proposition2}

Il en résulte que l'anneau $𝒟$ est local \emph{nœthérien}, régulier.

\begin{remarque2}
Notons qu'il est clair que $𝒟$ est « quasi-régulier ».
En effet, le gradué de $𝒫$ relativement à l'idéal $(x₁,…,x_n,y)$ est
une algèbre symétrique : pour chaque entier $r$,
le morphisme $A]\sous{t}[ [y] → 𝒫$ induit un isomorphisme
modulo $(x₁,…,x_n)^r$.
\end{remarque2}

\begin{démo}
Il suffit de montrer que le morphisme $𝒫 → \chap{A]\sous{t}[} [[y]]$ est
plat. D'après le critère de platitude rappelé
précédemment (\ref{18B-critère-platitude}), il suffit de montrer la platitude
sur l'ouvert $x₁\cdots x_n ≠ 0$.
En effet, le cas où seuls certains $x_i$
sont nuls se ramène à ce cas particulier :
en tensorisant par $A]\sous{t}[ [y]/x_i$ le morphisme
$𝒫 →  \chap{A]\sous{t}[} [[y]]$, on obtient une flèche
du même type définie par l'anneau $A/x_i$
de dimension $n-1$ et des séries qui coïncident
avec l'évaluation en $x_i=0$
des $f₁,…,\chap{f_i},…,f_n$.
Pour chaque $i$, $𝒟_i[x_i^{-1}]$ est une algèbre
de polynômes $A]\sous{t}[ [y,F_i][x_i^{-1}]$
de sorte que le morphisme dont on souhaite montrer la platitude
est
\[
A]\sous{t}[ [y,F₁,…,F_n] [\frac{1}{x₁ \cdots x_n}] →
\chap{A]\sous{t}[} [[y]] [\frac{1}{x₁ \cdots x_n}],
\]
\[
F_i↦ f_i.
\]
Il suffit de montrer que pour chaque sous-ensemble
\emph{fini} $𝒯$ des variables $\sous{t}$, le morphisme
\[
A]t ∈ 𝒯[ [y,F₁,…,F_n] [\frac{1}{x₁ \cdots x_n}]
→ \chap{A]t ∈ 𝒯[} [[y]] [\frac{1}{x₁ \cdots x_n}]
\]
est plat. Quitte à agrandir un tel ensemble $𝒯$,
on se ramène au cas où $𝒯$ est \emph{cofini}, de complémentaire
des variables $t_{1 α},…,t_{n α}$ pour un indice $α ≥ 1$
quelconque. Posant alors $R_α=A]t_{i,β≠α}[$
et $R′=A]\sous{t}[$, il suffit de montrer
que le morphisme $R_α[y,F₁,…,F_n]→ R′[[y]]$
est plat au-dessus de l'ouvert $x₁\cdots x_n ≠ 0$.
Un dernier dévissage nous ramène à montrer
la platitude du morphisme $R_α[y,F₁,…,F_n] →
R_α[y,t_{1 α},…,t_{n α}]$, au-dessus du même ouvert.
Ce dernier point résulte de la proposition
\ref{18B-platitude-série-formelle-universelle}.
\end{démo}

\begin{proposition2}
Le diviseur $C=V(y^ℓ+ F₁+\cdots+F_n)$ de $\Spec(𝒟)$
est \emph{régulier} hors du point fermé.
\end{proposition2}

Dans cet énoncé, on note abusivement $F_i$
l'image dans $𝒟$ de l'élément de $𝒟_i$ correspondant
à $f_i$ (cf. \ref{18B-dilatations-formelles}).

\begin{démo}
Il suffit de montrer que pour chaque sous-ensemble
strict $E$ de $[1,n]$, l'intersection schématique
de $C$ avec le sous-schéma
$X_E=\{x_i =0, i ∈ E ; x_i ≠ 0, i ∉ E\}$
de $\Spec(𝒟)$ est un diviseur régulier de $X_E$.
Si $E=∅$, cela résulte du fait que le schéma $X_∅$
est le localisé d'une algèbre de polynômes
en les $y,F₁,…,F_n$. Le cas général se ramène
aisément à ce cas particulier. (Remarquons que si
$x_i=0$, il en est de même de $F_i$.)
\end{démo}

\begin{proposition2}
L'image inverse de $C$ dans un localisé
strict de $\Spec(𝒟)$ est \emph{irréductible}.
\end{proposition2}

\begin{démo}
Même argument qu'en dimension $2$.
\end{démo}

\begin{corollaire2}
\label{18B-ouvert-dim-d-dimcoh-2d-1}
Pour tout entier $d ≥ 1$, il existe schéma nœthérien
strictement local régulier $X$ de dimension $d$ et un ouvert affine $U=D(f)$
de $ℓ$-dimension cohomologique $2d-1$.
\end{corollaire2}

\begin{démo}
La même démonstration qu'en dimension $2$
nous permet de minorer la dimension cohomologique
par $2d-1$. D'après \refext{XVIII-A}{18A-Gabber's bound1},
c'est une égalité.
\end{démo}

\section{Dimension cohomologique : majoration d'une « fibre de
Milnor générique »}
\label{18B-majoration-Milnor}

\subsection{Énoncé}

\begin{théorème2}
\label{18B-majoration Milnor générique}
Soit $R → R′$ un morphisme local essentiellement de type fini
d'anneaux nœthériens strictement locaux intègres.
Notons $K$ le corps des fractions de $R$.
Alors, pour tout nombre premier $ℓ$ inversible sur $R$,
on a la majoration
\[
\cd_ℓ(R′ ⊗_R K) ≤ \dim(R′),
\]
où le terme de gauche désigne la $ℓ$-dimension cohomologique
étale du spectre de l'anneau $R′ ⊗_R K$ et le terme
de droite désigne la dimension de Krull de $R′$.
\end{théorème2}

\subsubsection{}
\label{18B-essentiellement-tf}
Dans cet énoncé, l'hypothèse de finitude sur $f$
signifie que ce morphisme est localement de la forme
$\Spec(B) → \Spec(A)$, avec $B$ une colimite de $A$-algèbres
de type fini à morphismes de transition étales.

\begin{corollaire2}
Soit $R$ un anneau strictement local nœthérien intègre
de corps des fractions $K$ et soit $ℓ$ un nombre premier
inversible sur $R$. Alors, on a la majoration
\[
\cd_ℓ(K) ≤ \dim(R).
\]
\end{corollaire2}

\begin{remarque2}
Réciproquement, on peut montrer que si, sous les hypothèses du
corollaire, $U$ est un ouvert non vide strict de $\Spec(R)$, alors
$\mathrm{cd}_ℓ(U) ≥ \dim(R)$ et que, lorsque $R′ ⊗_R K ≠ 0$,
la majoration du théorème est une égalité.
Une façon de procéder, due à O. Gabber (non publié),
est d'utiliser une variante de la méthode (également due à O. Gabber) exposée
dans \cite[§6.1]{pdim@GO} et reposant sur une
« astuce quadratique ». Pour une autre méthode,
voir §\ref{18B-minoration-ouvert-spectre-épointé}, \emph{infra}.
Rappelons que la minoration « limite »
$\mathrm{cd}_ℓ(K) ≥ \dim(R)$ est élémentaire : on procède par
spécialisations successives en codimension $1$ (voir \sga{4}{X}{2.4}).
\end{remarque2}

\subsection{Démonstration}
\subsubsection{}On procède par récurrence sur $d′=\dim(R′)$ et
l'on se ramène au cas excellent.

\subsubsection{}
Notons $X=\Spec(R)$, $Y=\Spec(R′)$ et respectivement
 $X^⋆$ et $Y^⋆$ les spectres épointés.
Considérons l'ouvert $Y_⋆=Y×_X X^⋆$ de $Y^⋆$,
$V =\Spec(R′ ⊗_R K)$ la fibre générique de $Y → X$
et enfin $j$ le morphisme $V ↪ Y^⋆$.
Il résulte de l'hypothèse de récurrence que pour
chaque faisceau $ℱ$ de $𝐙/ℓ$-modules sur $V$,
le complexe $Φ_ℱ=𝖱 j_* ℱ$ appartient à $𝖣^{≤ \mathrm{cod}}(Y^⋆)$,
où $\mathrm{cod}$ est la fonction de perversité $y↦\dim 𝒪_{Y,y}$.
(Ceci est encore vrai avec $y↦d′-\dim \sur{\{y\}}$.)
On veut montrer que $H^r(V,ℱ)=H^r(Y^⋆,Φ_ℱ)$ est nul pour $r>d′$.
Fixons un tel $r$ et une classe $c ∈ H^r(Y^⋆,Φ_ℱ)$.

\subsubsection{}
On suppose $d ≥ 2$, et on choisit un système de paramètres $x₁,…,x_d$
pour l'anneau strictement local $R$.
Soit $Z=Y\}t₁,…,t_{d-1}\{$ le « gonflement étale »,
hensélisé strict de $𝐀^{d-1}_Y=Y[t₁,…,t_{d-1}]$ en un point générique
géométrique de la fibre spéciale sur $Y$.
L'« hyperplan » $H=V(t₁x₁+\cdots+t_{d-1}x_{d-1}+x_d)$ de $Z$
est de codimension $1$, essentiellement lisse au-dessus de $Y_⋆$.
Considérons le triangle distingué
\[
𝖱Γ_{H_⋆}(Z_⋆,Φ_ℱ) → 𝖱Γ(Z_⋆,Φ_ℱ) → 𝖱Γ(Z_⋆-H_⋆,Φ_ℱ) →,
\]
où l'on note $Z_⋆$ le produit fibré $Z ×_Y Y_⋆$ et, abusivement, $Φ_ℱ$
ses diverses images inverses. Soit $i$ l'immersion fermée $H_⋆ ↪ Z_⋆$,
où $H_⋆=H×_Z Z_⋆$.  Le morphisme $𝐙/ℓ → i^! 𝐙/ℓ(1)[2]$ de complexes
sur $H_⋆$ est un isomorphisme par pureté relative.  On montre par dévissage
qu'il en est de même de la flèche ${Φ_ℱ}_{|H_⋆} →
i^!{Φ_ℱ}_{|Z_⋆}(1)[2]$ obtenue par tensorisation à partir de la
précédente.  On utilise ici le fait que la restriction
de $Φ_ℱ$ à $Z_⋆$ provient de la base $Y_⋆$.  On en
tire le morceau de suite exacte :
\[
H^{r-2}(H_⋆,Φ_ℱ)(-1) → H^r(Z_⋆,Φ_ℱ) → H^r(Z_⋆-H_⋆,Φ_ℱ).
\]
Notons que $Z_⋆-H_⋆=Z-H$ car $H$ contient la fibre spéciale
de $Z → X$.
Soient $\chap{Y}$ le complété ($𝔪_{R′}$-adique) de $Y$
et $\gtilde{Z}$ un hensélisé strict du produit
fibré $Z×_Y \chap{Y}$.
Notons que le morphisme $\gtilde{Z} → Z$ est un morphisme local
entre schémas strictement locaux induisant
un isomorphisme sur la complétion le long
de la fibre spéciale sur $Y$.
Il résulte donc du théorème de comparaison
de Fujiwara-Gabber (\cite[6.6.4]{TUBULAR@Fujiwara}) que le morphisme
$H^r(Z_⋆,Φ_ℱ) → H^r(\gtilde{Z}_⋆,Φ_ℱ)$
est un \emph{isomorphisme} pour chaque $r$.
Il en est de même de $H^{r-2}(H_⋆,Φ_ℱ) →
H^{r-2}(\gtilde{H}_⋆, Φ_ℱ)$,
où $\gtilde{H}_⋆=H_⋆ ×_{Z_⋆} \gtilde{Z}_⋆$.
Le schéma $\gtilde{Z}_⋆-\gtilde{H}_⋆$ est un ouvert
affine, qui coïncide avec $\gtilde{Z}-\gtilde{H}$,
d'un schéma strictement local essentiellement de type fini
sur le schéma local nœthérien complet $\chap{Y}$.
Comme l'appartenance de $Φ_ℱ$ à $𝖣^{≤\mathrm{cod}}$
est préservée par complétion,
il résulte du théorème de Lefschetz affine
(\refext{XV}{15-prop-principale-locale}),
dans le cas excellent, que le groupe
de cohomologie $H^r(\gtilde{Z}_⋆-\gtilde{H}_⋆,Φ_ℱ)$
est nul pour chaque $r> \dim(Z)=\dim(Y)=d′$. En conséquence,
le morphisme $H^r(\gtilde{Z}_⋆,Φ) → H^r(\gtilde{Z}_⋆-\gtilde{H}_⋆,Φ)$
est nul pour les mêmes $r$. De ce fait, des théorèmes de comparaisons
susmentionnés et de la compatibilité du morphisme de Gysin
à la complétion, il résulte formellement que toute classe $c ∈ H^r(Z_⋆,Φ_ℱ)$
provient d'une classe dans $H^{r-2}(H_⋆,Φ_ℱ)(-1)$
et est donc tuée par restriction à $Z_⋆-H_⋆$.

\subsubsection{}
Il existe donc un voisinage étale $e:W → 𝐀^{d-1}_Y$ dont l'image
rencontre la fibre spéciale sur $Y$ tel que la classe $c ∈ H^r(Y^⋆,Φ_ℱ)$
soit tuée par restriction à $W-H_W$. Notons $k′$
le corps résiduel de $Y$, et $k$ celui de $X$.
L'ensemble $k^{d-1}$ est dense dans $𝐀^{d-1}_{k′}$, car $k$
est infini. Il en résulte qu'il existe une section $σ: Y →
𝐀^{d-1}_Y$, correspondant à des spécialisations des $t_i$
à valeurs dans $R$, telle que $W_σ$ ait une fibre spéciale
sur $Y$ non vide. Le schéma $Y$ étant strictement local, on relève
cette section en $Y → W_σ → W$. La classe de cohomologie $c$
est donc nulle sur $Y-H_Y$, où $H_Y$ est maintenant une hypersurface
d'équation $x_d+t₁x₁+\cdots+t_{d-1}x_{d-1}$ à coefficients $t_i$ dans $R$.
Cet ouvert affine $Y-H_Y$ contient la fibre générique $V=Y ⊗_R K$
car l'élément $x_d+t₁x₁+\cdots+t_{d-1}x_{d-1} ∈ R$ est non nul,
les $x_i$ constituant un système de paramètres de $R$.
Finalement la restriction de $c ∈ H^r(Y^⋆,Φ_ℱ)$ à $H^r(V,Φ_ℱ)=H^r(V,ℱ)$,
qui est la classe dont on est parti, est nulle. CQFD.

\section{Majoration : amélioration}
\label{18B-majoration-générale}

\subsection{Énoncé}

\subsubsection{}Soit $f:Y → X$ un morphisme entre espaces topologiques
sobres non vides.
On note
\[
\mathrm{dim.cat}(f)=\sup \{n ∈ 𝐍 : ∃ y₀ ⤳  y₁ ⤳  \cdots ⤳ y_n,
f(y₀) ≠ f(y₁) ≠ \cdots ≠ f(y_n) \} ∈ 𝐍 ∪ \{∞\}
\]
la \emph{dimension caténaire} de $f$, où chaque $y_i ⤳ y_{i+1}$ est une spécialisation.

\subsubsection{}
Par construction, $\mathrm{dim.cat}(f:Y → X)$ est majorée
par les dimensions de $X$ et de $Y$ avec égalité par exemple
lorsque $f$ est l'identité.
Plus généralement, lorsque $f$ est un morphisme \emph{générisant} (\ega{I′}{3.9.1})
— comme c'est le cas d'un morphisme plat de schémas —
la dimension caténaire coïncide avec la dimension de
l'image.

\begin{remarque2}
Si $f$ est un morphisme dominant essentiellement de type
fini (au sens du \ref{18B-essentiellement-tf}) entre schémas nœthériens intègres, on peut montrer
que la dimension caténaire de $f$
est la dimension de l'image d'une platification
de $f$.
\end{remarque2}

\begin{théorème2}
Soit $f:Y → X$ un morphisme essentiellement de type fini
entre schémas nœthériens strictement locaux et soit $V$ un
ouvert affine de $Y$. Alors, pour tout nombre premier $ℓ$
inversible sur $X$, on a la majoration
\[ \mathrm{cd}_ℓ(V) ≤
\dim(Y)+\max\big(0,\mathrm{dim.cat}(f)-1\big).  \] En
particulier, si $\dim(X) ≥ 1$, $\mathrm{cd}_ℓ(V) ≤
\dim(Y)+\dim(X)-1$.
\end{théorème2}

\begin{corollaire2}
Soit $f:Y → X$ un morphisme essentiellement de type fini
entre schémas nœthériens strictement locaux, où $\dim(X) ≥ 1$,
et soit $V$ un ouvert affine de $Y$. Alors, pour tout nombre premier $ℓ$
inversible sur $X$, on a la majoration
\[ \mathrm{cd}_ℓ(V) ≤
\dim(Y)+\dim(X)-1.\]
\end{corollaire2}

\begin{corollaire2}
\label{18B-n-entre-d-et-2d-1}
Soit $d ≥ 1$ un entier et soit $n$ un entier dans
l'intervalle fermé $[d,2d-1]$. Il existe un schéma
nœthérien strictement local $X$, régulier de dimension $d$,
et un ouvert affine $U$ de ce schéma tel pour tout nombre
premier $ℓ$ inversible sur $X$ on ait l'égalité
\[
\mathrm{cd}_ℓ(U)=n.
\]
\end{corollaire2}

\begin{démo}
Il suffit de montrer que pour tout entier $d ≥ 1$, et tout entier $r ≥ 0$,
il existe un schéma nœthérien strictement local régulier $Y$ de
dimension $d+r$ et un ouvert affine $V$ de $Y$
de $ℓ$-dimension cohomologique égale
à $2d+r -1$. Soient $X$ et $f$ comme en
\ref{18B-ouvert-dim-d-dimcoh-2d-1} :
l'ouvert affine $U=X[f^{-1}]$ est
de dimension $d$, $ℓ$-dimension cohomologique $δ=2d-1$.
Plus précisément, il résulte de la démonstration
qu'il existe une classe non nulle dans $H^δ(U,𝐙/ℓ)$.
Considérons maintenant $Y=X[T]_{(0)}$ un hensélisé stricte
de la droite affine sur $X$ en l'origine de la fibre
spéciale, $g=fT ∈ Γ(Y,𝒪_Y)$ et $V$ l'ouvert affine
$Y[g^{-1}]$. On a $\dim(Y)=d+1$. Par pureté cohomologique,
on vérifie immédiatement que le groupe de cohomologie $H^{d+1}(V,𝐙/ℓ)$
est également non nul. Par récurrence, on obtient une
paire $(Y,V)$ comme ci-dessus telle que $\mathrm{cd}_ℓ(V) ≥ 2d+r-1$.
D'après le corollaire précédent, on a également la majoration
$\mathrm{cd}_ℓ(V) ≤ 2d+r-1$, d'où l'égalité.
\end{démo}

\subsection{Démonstration}

\subsubsection{}On procède par récurrence sur la dimension caténaire de $f$.

\subsubsection{$\mathrm{dim.cat}(f)=0$} Cette égalité
se produit si et seulement si $Y$ est contenu dans la fibre spéciale. Le théorème
est donc connu dans ce cas : on est sur un corps
donc dans une situation excellente.

\subsubsection{$\mathrm{dim.cat}(f)=1$}
On peut supposer les schémas $X$ et $Y$ réduits.
Quitte à procéder par récurrence sur la dimension de $Y$, on peut également
supposer $Y$ irréductible : si $Y$ est la réunion
de deux fermés stricts $Y₁$ et $Y₂$, considérer
par exemple le morphisme $π:Y₁ ⨿ Y₂ → Y$ et
la suite exacte $0 → ℱ → π^* ℱ → i_*ℋ → 0$, où $i$
est l'immersion fermée $Y₁ ∩ Y₂ ↪ Y$ et $ℋ$ un faisceau sur
cette intersection. Quitte à remplacer $X$
par l'adhérence de l'image de $f$, on peut également supposer
la base intègre et $f$ dominant.
Soit $η$ (resp. $s$) le point générique (resp. fermé) de $X$
et $η′$ (resp. $s′$) le point générique (resp. fermé) de $Y$.
Comme tout point $y$ de $Y$ s'insère dans
une suite de spécialisations $η′ ⤳ y ⤳ s′$, d'image
$η ⤳ f(y) ⤳ s$, il résulte de l'hypothèse $\mathrm{dim.cat}(f)=1$
que $f(y)=η$ ou $f(y)=s$. Soient $\chap{X}$ le complété de $X$ et $\gtilde{Y}$ un hensélisé
strict du produit fibré $Y×_X \chap{X}$. C'est un schéma
strictement local de dimension $\dim(Y)$ et
\emph{excellent} car essentiellement de type fini sur le schéma
local nœthérien complet — donc excellent — $\chap{X}$.
On note $\gtilde{V}$ l'ouvert $V×_Y \gtilde{Y}$.
Il résulte de \sga{4½}{[Finitude]}{1.9} qu'au-dessus de $η$,
et donc au-dessus de $X-\{s\}$, la formation des images directes par $j:V ↪ Y$ commute au changement
de base $\gtilde{Y} → Y$.
En d'autres termes, si $V′=V ∪ (Y-Y_s)$ et $j′$ désigne
l'immersion intermédiaire $V ↪ V′$,
la formation de $𝖱 j′_*$ commute à $\gtilde{Y} → Y$.
Il en est de même pour $j″: V′ → Y$ d'après le théorème
de comparaison à la complétion de Fujiwara-Gabber
(\cite[6.6.4]{TUBULAR@Fujiwara}).
On utilise ici le fait que si $F$ est un fermé de $Y$
inclus dans la fibre spéciale alors les complétés
de $Y$ le long de $F$ et celui de l'hensélisé
de $\gtilde{Y}$ (en le point correspondant au point
fermé de $Y$) sont naturellement isomorphes.
Finalement, le foncteur $𝖱Γ(V)=𝖱Γ(Y) ∘ 𝖱j_*$ s'identifie
au foncteur $𝖱Γ(\gtilde{Y}) ∘ 𝖱 \gtilde{j}_*=𝖱Γ(\gtilde{V})$,
appliqué à l'image inverse.
Ainsi on a l'inégalité $\mathrm{cd}_ℓ(V) ≤ \mathrm{cd}_ℓ(\gtilde{V})$.
Le schéma $\gtilde{Y}$ est quasi-excellent, le terme de
droite est donc justiciable du théorème de Lefschetz affine.
L'inégalité $\mathrm{cd}_ℓ(V) ≤ \dim(Y)$ résulte alors
de l'égalité $\dim(Y)=\dim(Y′)$.

\begin{remarque2}
Lorsque $\dim(R)=1$, on a vu en \refext{XIII}{13-compléments annulation}
que le théorème peut également se démontrer
par normalisation.
\end{remarque2}

\subsubsection{$\mathrm{dim.cat}(f)>1$}
Notons à nouveau $j$ l'immersion ouverte $V ↪ Y$.
Par restriction à la fibre générique, on a un isomorphisme
$𝖱Γ(V_η,ℱ)=𝖱Γ(Y_η,(R j_* ℱ)_{|Y_η})$.
Si $y$ est un point géométrique de $Y$ localisé en $Y_η$,
la fibre $(R^q j_* ℱ)_y$ est nulle dès lors que $q>\dim(𝒪_{Y_η,y})$.
Cela résulte par passage à la limite du théorème d'Artin
pour les schémas affines de type fini sur un corps et du fait trivial
que $Y_{(y)} → Y$ se factorise à travers $Y_η$.
Soient $q ≥ 0$ un entier, $𝒢$ un sous-faisceau constructible de $(R^q j_* ℱ)_{|Y_η}$
et $S$ l'adhérence dans $Y$ de son support.
D'après ce qui précède, on a la majoration
$\codim(S_η,Y_η) ≥  q$. Il en résulte que $\codim(S,Y) ≥ q$,
et ce sans hypothèse de caténarité sur les schémas.
De la suite spectrale de composition des foncteurs
et du théorème \ref{18B-majoration Milnor générique},
on déduit que le groupe de cohomologie $H^n(Y_η,𝒢)$ est nul lorsque
$n>\dim(Y)$.

Considérons la flèche d'adjonction $ℱ → k_*k^* ℱ$,
où $k$ est l'immersion $V_η ↪ V$, et $𝒦$ son noyau.
Par construction, la restriction de $𝒦$ à $V_η$ est nulle.
La dimension de l'adhérence du support de $𝒦$ est donc
au plus $\dim(Y)-1$. Il résulte donc de l'hypothèse de récurrence
que le résultat d'annulation désiré est connu pour $𝒦$.
Procédant de même pour le conoyau de l'adjonction précédente,
on se ramène à démontrer l'annulation du groupe
$H^p(V,R⁰k_*k^* ℱ)$ pour $p ≥ \dim(Y)+\mathrm{dim.cat}(f) (>\dim(Y))$.
Compte tenu du résultat d'annulation précédemment établi
pour $𝖱Γ(V_η,ℱ)=𝖱Γ(V,R k_* k^* ℱ)$ et de la suite spectrale de Leray
$E₂^{p,q}=H^p(V,R^q k_* k^* ℱ) ⇒ H^{p+q}(V_η,ℱ)$,
il suffit de montrer que pour chaque tel $q$,
les groupes $H^{p-q-1}(V,R^q k_*k^* ℱ)$ sont nuls pour $q>0$.
Fixons $q$. Soit $y$ un point géométrique de $Y$ tel
que la fibre de $R^q k_*k^* ℱ$ en $y$ soit non nulle et $x$ le point
image de $Y$ dans $X$. Le schéma $η×_X X_{(x)}$ se décompose
en un coproduit de spectres de corps $η_α$ ; de même, le produit fibré
$Y_{(y)}×_X η$, dont on considère la cohomologie, est isomorphe au coproduit
des $Y_{(y)}×_{X_{(x)}} η_α$.
D'après \emph{op. cit.} (\ref{18B-majoration Milnor générique}),
ces derniers n'ont de cohomologie qu'en degré $q ≤ \dim Y_{(y)} ≤
\dim(Y)-\dim \sur{\{y\}}$. Il en résulte que la dimension
du support de chacun des sous-faisceaux constructibles de $R^q k_*k^*
ℱ$ est au plus $\dim(Y)-q$. De plus, la dimension caténaire du
morphisme $f$ restreint à un tel support est au plus
$\mathrm{dim.cat}(f)-1$. Il résulte donc de l'hypothèse
de récurrence que les groupes
$H^{p-q-1}(V,R^q k_*k^* ℱ)$ sont nuls lorsque
$p-q-1 ≥ (\dim(Y)-q)+(\mathrm{dim.cat.}(f)-1)$. CQFD.

\section{Dimension cohomologique d'un ouvert du spectre épointé : minoration}
\label{18B-minoration-ouvert-spectre-épointé}

\begin{théorème}
Soient $X$ un schéma intègre strictement local nœthérien de
dimension $d$ et $Ω$ un ouvert non vide du spectre épointé. Alors, pour tout nombre
entier $n>1$ inversible sur $X$, on a
\[
\mathrm{cd}_n(Ω) ≥ d.
\]
\end{théorème}

La démonstration occupe les paragraphes suivants.

\subsection{Construction combinatoire locale}
\label{18B-construction-combinatoire-locale}
\subsubsection{Notations}
Soit $X$ un schéma strictement local nœthérien régulier de
dimension $d ≥ 2$ et soit $t₁,…,t_{d-1},t_d$ un système régulier de
paramètres. Pour des raisons qui apparaîtront ultérieurement, on
note également $π$ l'élément $t_d$. Pour chaque $1 ≤ i ≤ d-1$, on note $H_i$ le
diviseur régulier $V(t_i)$ ; pour $i=d$, on pose
$H_d=V(t₁+\cdots+t_{d-1}-π)$. Enfin on note $U$ l'ouvert
affine $X[π^{-1}]$, $k$ l'immersion ouverte $U ↪ X$
et $j$ l'immersion ouverte $U-⋃_{i=1}^d H_i ↪ U$.
On fixe un entier $n$ inversible sur $X$ et on pose $Λ=𝐙/n$.

\subsubsection{}
\label{18B-cohomologie-via-pureté}
Soit $P$ une partie de $\{1,…,d\}$.
Notons $H_P$ l'intersection $⋂_{p ∈ P} H_p$, et désignons
par $H′_P$ l'intersection $H_P ∩ U$, ouverte dans $X$ et $H_P$
mais fermée dans $U$, et $k_P$ l'immersion ouverte $H′_P ↪ H_P$.
Pour chaque entier $q$, le groupe de cohomologie $H^q(U,Λ_{H′_P})$ est isomorphe
au groupe $H^q(H_P,Λ_{H′_P})$. Comme $H_P-H′_P$ est le diviseur
régulier défini par $π$ dans $H_P$, il résulte de la pureté
cohomologique que $H^q(U,Λ_{H′_P})$ est nul pour $q>1$,
isomorphe à $Λ$ pour $q=0$ et de rang $1$, engendré par
la classe de Kummer de $π$ pour $q=1$. Ceci vaut
également pour $P=∅$, avec la convention évidente
que $H_∅=X$ et $H′_∅=U$.

\subsubsection{}
\label{18B-qis}
Considérons maintenant le quasi-isomorphisme
\[
j_! Λ ⥲ \big( Λ → ⨁_{1 ≤ i ≤ d} Λ_{H′_i} → \cdots → ⨁_{|P|=d-1} Λ_{H′_P}
→ 0\big)
\]
de faisceaux de $Λ$-modules sur $U$,
où le premier terme du complexe de droite est placé en degré $0$.
(On utilise ici le fait que $H′_{P}=∅$ si $|P|=d$
car $H_P$ est alors le point fermé de $X$.)
Les différentielles sont des sommes, avec des signes, de flèches de restriction.
À cette résolution est associée — via la filtration « stupide » —
la suite spectrale
\[
E₁^{p,q}=H^q(U, ⨁_{|P|=p} Λ_{H′_P}) ⇒ H^{p+q}(U,j_! Λ).
\]
D'après les observations du paragraphe précédent,
toute classe de $E₁^{d-1,1}=⨁_{|P|=d-1} H¹(U,Λ_{H′_P})$
n'appartenant pas à l'image de $E₁^{d-2,1}$
survit dans l'aboutissement $H^d(U,j_! Λ)$.

\subsubsection{}
\label{18B-classe-locale-non-nulle}
De même que le faisceau $j_! Λ$ est isomorphe au produit tensoriel
\[
(U-H′₁ ↪ U)_! Λ ⊗ \cdots ⊗ (U - H′_d ↪ U)_! Λ),\]
le complexe quasi-isomorphe à $j_! Λ$ ci-dessus
est isomorphe au produit tensoriel des complexes $(Λ → Λ_{H′_i})$, $1 ≤ i ≤ d$, où la flèche
est l'unité de l'adjonction, isomorphes respectivement aux $(U-H′_i ↪ U)_! Λ$.
De cette observation, jointe à (\ref{18B-cohomologie-via-pureté}),
on en déduit que le complexe $E₁^{∙,1}$
est isomorphe à la troncation naïve $σ_{≤d-1}\big((Λ
\dessusdessous{\Id}{→} Λ)^{⊗d}\big)$
obtenue en remplaçant le $d$-ième terme, isomorphe à $Λ$, par zéro.
Il est bien connu que ce complexe (« de Koszul ») produit tensoriel
est acyclique (avant troncation), cf. par exemple \ega{III₁}{§1.1}.
(L'exactitude résulte également du quasi-isomorphisme ci-dessus, appliqué
à d'autres fermés.)
En particulier, l'image de la différentielle
$E₁^{d-2,1} → E₁^{d-1,1}$ est naturellement le noyau d'une forme linéaire
non nulle (explicite) sur $E₁^{d-1,1}$ et n'est donc pas $E₁^{d-1,1}$ tout entier.
Il existe donc des sommes directes de classes de Kummer de $π$
qui survivent dans $H^d(U,j_! Λ)$.

\subsection{Éclatement et normalisation partielle}

\subsubsection{}
\label{18B-stratégie-éclatement-normalisation-partielle}
Soit maintenant $X=\Spec(R)$ un schéma strictement local nœthérien intègre
de dimension $d ≥ 2$, de point fermé $x$, et soit $Ω$ un ouvert non vide
strict de $X$. Nous allons montrer qu'après éclatement
et « normalisation partielle » l'ouvert $Ω$ est
— localement et « modulo des nilpotents » —
un schéma régulier du type du schéma $U$ considéré ci-dessus.
Ceci permet de produire une classe de cohomologie non nulle de degré $d$
sur $Ω$.

\subsubsection{}Soit $Y=\Ecl_x(X)$. Notons $j$ l'immersion ouverte
de $Y-Y_x$ dans $Y$.
On désigne par $\chap{X}$ le complété du schéma local $X$,
par $Y′$ le produit fibré $Y×_X \chap{X}$ et par $Y′_{\red}$ la réduction de $Y′$.
Notons que le schéma $Y′$ est \emph{excellent} car $\chap{X}$
l'est.

\subsubsection{}Soit $𝒪_Y^{\japonaismath{の}}$ la normalisation de $𝒪_Y$
dans $j_* 𝒪_{Y-Y_x}$. On définit de même
$𝒪_{Y′}^{\japonaismath{の}}$ et $𝒪_{Y′_\red}^{\japonaismath{の}}$.
La $𝒪_Y$-algèbre $𝒪_Y^{\japonaismath{の}}$ est colimite (filtrante) de ses
sous-$𝒪_Y$-algèbres de type fini $𝒜_λ$.

\begin{proposition2}
Le foncteur envoyant une sous-$𝒪_Y$-algèbre $ℬ$
de $j_* 𝒪_{Y-Y_x}$ sur l'image de $(Y′ → Y)^*ℬ$ par l'adjonction
$(Y′ → Y)^*j_*𝒪_{Y-Y_x} → j′_* 𝒪_{Y′ -Y_x}$
induit une \emph{bijection} entre
les sous-$𝒪_Y$-algèbres finies de $j_* 𝒪_{Y-Y_x}$
et les sous-$𝒪_{Y′}$-algèbres finies de $j′_* 𝒪_{Y′-Y_x}$.
De plus, les algèbres $𝒪_Y^{\japonaismath{の}}$ et
$𝒪_{Y′}^{\japonaismath{の}}$ se correspondent par ce foncteur.
\end{proposition2}

\begin{démo}
On se ramène au cas où $Y=\Spec(A)$ et $Y_x=V(t)$.
Il suffit de montrer que si $A$ est une $R$-algèbre intègre
et $A′=A ⊗_R \chap{R}$ alors le morphisme
$A[t^{-1}]/A →  A′[t^{-1}]/A′$ est un isomorphisme
et que les normalisations se correspondent.
Notons que les anneaux $A$ et $A′$ ont même complétion $t$-adique.
Le premier point résulte alors du fait que si $M$ est un $A$-module
dont chaque élément est tué par une puissance de $t$, on a $M ⥲ M ⊗_A A′$.
Enfin, soit $(f′/t^r)^d + a′₁ (f′/t^r)^{d-1}+ \cdots +a′_d =0$
une relation de dépendance intégrale où $f′$ et les $a′_i$
appartiennent à $A′$.
Soit $N$ un entier assez grand. Écrivons $f′ = f+  t^N g′$,
$a′_i=a_i+ t^N b′_i$ où $f$ et les $a_i$ appartiennent à $A$.
La relation précédente devient
$(f/t^r)^d + a₁ (f/t^r)^{d-1}+ \cdots +a_d ∈ A′ ∩ A[t^{-1}]=A$.
Il en résulte que l'élément $f′/t^r=f/t^r + t^{N-r}g′$
est, modulo un élément de $A′$, dans l'image
de $A^{\!\japonaismath{の}}$.
\end{démo}

\subsubsection{}Par excellence, l'algèbre $𝒪_{Y′_\red}^{\japonaismath{の}}$
est finie sur $𝒪_{Y′_{\red }}$.
L'anneau $𝒪_{Y′} ^{\japonaismath{の}}$ en est l'« image inverse »
par la surjection naturelle. 
De ces observations et de la proposition précédente, on déduit
qu'il existe un indice $λ$ tel que, si $Z=\Spec(𝒜_λ)$
et $Z′= Z_{\chap{X}}$, alors $Z′_{\red}$ est intégralement clos
dans $Z′_{\red}-E′_{\red}$. Notons que $Z$ et $Y$ sont isomorphes hors
de $E$.

\subsubsection{}Soit $E$ une composante de dimension $d-1$ de $Y_x$
et soit $e$ un point maximal de $E′_{\red}$ dans $Z′_\red$.
Le localisé en $e$ est un anneau de valuation discrète : c'est un
anneau local réduit de dimension $1$, intégralement clos dans le
complémentaire du point fermé. Par excellence de $Z′$
il existe un ouvert dense de $E′_\red$ le long duquel $Z′_\red$
est régulier. On peut également supposer que $E′_\red$
est régulier sur cet ouvert. (Pour ce dernier
point il suffit de constater que $E′_\red$ est de type fini sur un corps.)
Soit $U′_\red$ un ouvert de $Z′_\red$
induisant l'ouvert de $E′_\red$ ci-dessus et $U$ un
ouvert de $Z$ induisant l'ouvert correspondant de $E$.
(Le morphisme $Z′ → Z$ est un isomorphisme sur $E$.)
On a $U′ ⊆ U×_Z Z′$. 
Ci-dessous, on s'autorise à rétrécir
les ouverts $U$ et $U′$, sous réserve
qu'ils contiennent tous les points maximaux de $E$.
On suppose de plus que $U ∩ Y_x = U ∩ E$.

\subsubsection{}On note $t$ une équation de $E$ dans $U$ et $π$
une équation de $E′_{\red}$ dans $U′_{\red}$
de sorte qu'il existe une unité $u$ et un entier $e$
tels que l'on ait l'égalité $t=u × π^e$ sur $U′_{\red}$.
L'existence d'un relèvement montre que l'on peut supposer
l'équation $π$ définie sur $U′$. Vérifions que l'\emph{on peut également
supposer $π$ définie sur $U$}. Les schémas $Z′$ et $Z$ ayant même
complétion $t$-adique, il suffit d'observer que si $a$ est une
fonction sur $U′$, on a l'égalité d'idéaux
$(π)=(π+at²)$, du moins lorsque
$1+ uaπ^{2e-1} ∈ 𝐆_m(U′)$, ce que l'on peut supposer quitte à
restreindre $U′$.


\subsubsection{}Soit $Ω$ un ouvert non vide de $X-\{x\}$
(cf. \ref{18B-stratégie-éclatement-normalisation-partielle}).
On note également $Ω$ ses images inverses dans $Y$ et $Z$ ;
elles lui sont isomorphes. Sur un voisinage ouvert des points maximaux
de $E$, l'ouvert $Ω$ coïncide avec le complémentaire de $E$.
Génériquement sur $E$, on a donc $Ω=Z[π^{-1}]$.
Soit $z$ un point fermé de $Z$ appartenant à un tel ouvert
ainsi qu'à l'ouvert $U$. Soient $t₁,…,t_{d-1}$ des fonctions de $𝒪_{Z ′_{\red},z}$
constituant, avec $π$, un système régulier de paramètres. On peut
supposer qu'elles s'étendent à $U′_\red$.
Utilisant à nouveau le fait que le morphisme $Z′ → Z$ est
un isomorphisme au-dessus de $E$, on peut également
supposer qu'elles proviennent de $U$, quitte
à les changer modulo $π$.
Pour chaque $i<d$, considérons l'adhérence schématique $H_i ⊆ Z$
de l'hypersurface $V(t_i)$ dans $U$  et $H_d$
l'adhérence schématique de $V(t₁+\cdots+t_{d-1}-π)$.
On note $H_i′=H_i ∩ Ω$.

\subsubsection{Stratégie}
On va construire une classe non nulle dans
le groupe de cohomologie $H^d(Ω,j_! Λ_{Ω-⋃₁^d H′_i})$,
où $j$ est l'immersion ouverte $Ω-⋃₁^d H′_i ↪ Ω$
et $Λ=𝐙/n$ avec $n$ inversible sur $X$.
\emph{Localement}, ces groupes de cohomologie sont invariants
par passage à la complétion de la base $X$
(et bien sûr au schéma réduit) de sorte que l'on va pouvoir utiliser les
calculs de \ref{18B-construction-combinatoire-locale}.
Il faut cependant prendre garde ici au fait que
l'intersection $⋂₁^d H′_i$ n'est pas nécessairement vide,
contrairement au cas local précédemment étudié : l'analogue
du complexe \ref{18B-qis} a donc un terme de plus (en degré $d$).
Malgré tout, on va relever à $Ω$ une classe de degré $d$
« locale » — c'est-à-dire du schéma $Ω×_Z Z_{(\sur{z})}$
(ou plutôt l'analogue sur $Z′_\red$) — à coefficients dans $j_!Λ$.

\subsubsection{}Considérons à nouveau la suite spectrale du
\ref{18B-qis} :
\[
E₁^{p,q}=H^q(Ω, Λ_p) ⇒ H^{p+q}(Ω,j_! Λ),
\]
où l'on note $Λ_p$ la somme directe des $Λ_{H′_P}$ avec $|P|=p$.
Pour chaque $P ⊆ [1,d]$ de cardinal $d-1$, l'intersection $H_P$
des hypersurfaces correspondantes de $Z$ est propre sur $X$.
Par construction, elle est aussi \emph{quasi-finie} au voisinage
du point fermé $z$. Il en résulte par le théorème de changement de
base propre pour le $π₀$ (ou le théorème principal de Zariski)
— d'après lequel sa décomposition en
composantes connexes se lit sur la fibre spéciale —
que chaque $H_P$ se décompose en le coproduit d'un schéma local
fini sur $X$ et d'un schéma ne rencontrant pas $z$.
Ainsi, $H^1(Ω,Λ_{d-1})$ a un facteur direct
isomorphe à $H¹(Z_{(z)}[π^{-1}],Λ_{d-1})$
et, par comparaison à la complétion, à $H¹({Z′_\red}_{(z)}[π^{-1}],Λ_{d-1})$.
De plus, la différentielle $E₁^{d-1,1} → E₁^{d,1}$
envoie ce facteur sur $0$ dans $H¹(Ω,Λ_d)$ car
l'intersection de $H_{[1,d]}$ avec $Ω$ est vide au voisinage de $z$.
En effet, si $t₁=t₂=\cdots=t_{d-1}=t₁+\cdots+t_{d-1}-π=0$
alors $π=0$ ; d'autre part, au voisinage de $z$, $Ω=Z[π^{-1}]$.
Ainsi, toute classe de cohomologie du facteur
direct « local » induit une classe dans $H^d(Ω,j_!Λ)$
qui relève la classe correspondante dans
$H^d({Z′_\red}_{(z)}[π^{-1}],j_!Λ)$.
On a vu (\ref{18B-classe-locale-non-nulle}) qu'il existe de telles
classes non nulles.
En conséquence, $H^d(Ω,j_! Λ) ≠ 0$ et
finalement $\cd_n(Ω) ≥ d$. CQFD.

\ifx\danslelivre\undefined
\bibliography{../fichiers-configuration-latex/bibliographie}
\bibliographystyle{../fichiers-configuration-latex/style-GTG}
\end{document}
\fi
